\documentclass{article}
\usepackage[margin=25mm]{geometry}
\usepackage{amsmath,amssymb}
\usepackage{xurl}
\usepackage[hidelinks]{hyperref}
\setlength{\emergencystretch}{2em}
% Shared commands for notes in docs/paper-gaps/.

\DeclareUnicodeCharacter{2080}{\ifmmode{}_{0}\else\textsubscript{0}\fi}
\DeclareUnicodeCharacter{2081}{\ifmmode{}_{1}\else\textsubscript{1}\fi}
\DeclareUnicodeCharacter{2082}{\ifmmode{}_{2}\else\textsubscript{2}\fi}
\DeclareUnicodeCharacter{2099}{\ifmmode{}_{n}\else\textsubscript{n}\fi}

\newcommand{\Lean}{\textsc{Lean}}
\ExplSyntaxOn
\NewDocumentCommand{\leanid}{m}{
  \ifmmode
    \textsf{\detokenize{#1}}
  \else
    \group_begin:
    \urlstyle{sf}
    \tl_set:Nn \l_tmpa_tl {#1}
    \tl_replace_all:Nnn \l_tmpa_tl {\_} {_}
    \exp_args:NV \path \l_tmpa_tl
    \group_end:
  \fi
}
\ExplSyntaxOff
\providecommand{\tr}{\operatorname{tr}}
\newcommand{\tnleanIssueBase}{https://github.com/LionSR/TNLean/issues/}
\newcommand{\tnleanPrBase}{https://github.com/LionSR/TNLean/pull/}
\newcommand{\tnleanDocBase}{https://sirui-lu.com/TNLean/docs/}
\newcommand{\ghissue}[1]{\href{\tnleanIssueBase#1}{GitHub issue \##1}}
\newcommand{\ghpr}[1]{\href{\tnleanPrBase#1}{PR \##1}}
\newcommand{\doclink}[1]{\href{\tnleanDocBase#1.html}{\path{#1}}}

% Dirac notation and the identity operator, as in blueprint/src/macros/common.tex.
% \providecommand keeps note-local definitions authoritative.
\usepackage{dsfont}
\providecommand{\ket}[1]{|#1\rangle}
\providecommand{\bra}[1]{\langle #1|}
\providecommand{\braket}[2]{\langle #1 | #2 \rangle}
\providecommand{\ketbra}[2]{|#1\rangle\!\langle #2|}
\providecommand{\Id}{\mathds{1}}
\providecommand{\one}{\mathds{1}}
\providecommand{\C}{\mathbb{C}}
\providecommand{\MN}[1]{M_{#1}(\C)}
\providecommand{\MD}{M_D(\C)}

% Verdict marker: \gapnote{<kind>}{<status>}. Kind is the policy
% classification (clarification, local-correction, scope-restriction,
% unfaithful, false-source, open-gap); status is open, wip (elimination
% underway), resolved, or historical. Machine-read by the site index;
% prints a verdict line.
\newcommand{\gapnote}[2]{\par\noindent\textbf{Verdict:} #1 (#2).\par}

\title{Energy centers and the central soft regional operator}
\author{}
\date{2026-11-03}
\begin{document}
\maketitle
\gapnote{scope-restriction}{open}

\section{The source interfaces}

Definition \texttt{def:soft-trunc} of
\href{https://arxiv.org/abs/1602.08828v2}{Arad--Landau--Vazirani--Vidick},
source lines 1178--1185, assumes a one-sided estimate
$e_J-10\le p_J\le e_J$ for the minimum $e_J$ of $H_J$, and defines
$p_JI+h_{m,t}(H_J-p_JI)$. The efficient-generation theorem
\texttt{t:aisefficient}, lines 1360--1369, instead permits an input
$|p_M-e_M|\le10$. The two interfaces are different: the latter condition
alone does not make $H_M-p_MI$ positive. An application of the one-sided
functional-calculus bounds must first derive a suitable one-sided center.

The regions in \texttt{eq:defjgapped}, lines 1380--1384, omit collars of
width $\ell$ at the boundaries of the parent region. Thus a rough
estimate for the parent minimum is not, without further argument, a
width-ten estimate for the middle minimum. One possible corrected route
uses the actual minima $c_L,c_R$ of the two short collars. The elementary
three-region comparison
$e_{\mathrm{mid}}+c_L+c_R\le e_M\le e_{\mathrm{mid}}+c_L+c_R+2$
would turn a parent estimate $p_M\in[e_M-8,e_M]$ into
$p_M-c_L-c_R-2\in[e_{\mathrm{mid}}-10,e_{\mathrm{mid}}]$.
This conversion, including the construction of the collar minima, is
separate from the results recorded here.

The raw soft replacement retains $p_JI$. Its norm can therefore contain
an extensive scalar energy. The uniform norm conclusion in
\texttt{t:aisefficient} and the norm assertion following
\texttt{eq:defjgapped} require an explicit scalar-centering convention;
they are not established here for the literal uncentered operator.
The correction to the harmonic cap, from $t\log m$ to $t(1+\log m)$,
is documented in the
\href{https://sirui-lu.com/QICLean/paper-gaps/alvv16_soft_truncation_log_bound.pdf}{QICLean soft-truncation note}.

\section{An actual centered operator and its energy}

Let $0\le h_i\le I$ be nearest-neighbour interactions on a finite chain
of positive local dimension. Fix $l\le u\le v\le r$ and let $P$ and $M$
be the sums internal to $[l,r)$ and $[u,v)$, respectively. Assume that
both collars have at most $2\ell$ sites. The formalization constructs
attained minima $e_P,e_M$ before every $p,t,m$. For
$e_M-10\le p\le e_M$, $t>0$ and $m\ge2$, it uses the literal operator
\begin{align}
 R=P-M+pI+h_{m,t}(M-pI).
\end{align}
It constructs an attained minimum $e_{\mathrm s}$ and an actual positive
matrix $K$ on the Hilbert space of $[l,r)$ such that
\begin{align}
 R&=K\otimes I+e_{\mathrm s}I,&
 \|K\|&\le t(1+\log m)+4\ell+14.
\end{align}
The identity extension is understood through the standard regional
coordinate identification. The operator $R$ is defined before
$e_{\mathrm s}$: its construction does not require knowing its true
minimum. If $T,t\ge48$, put
\begin{align}
 \rho=32e^{-(T-48)/8}+\frac{t}{m}\left(\frac{T}{t}\right)^m.
\end{align}
Then $e_{\mathrm s}\le e_P\le e_{\mathrm s}+\rho$. For every vector $g$
in the full ground space of the original chain,
\begin{align}
 \langle g,(K\otimes I)g\rangle\le(2+\rho)\|g\|^2.
\end{align}
No gap or uniqueness assumption is used. The matrices, their minima and
$K$ are chosen before the auxiliary cutoff $T$ and before any candidate
subspace. The formal statement is
\leanid{TNLean.MPS.AreaLaw.ALVV.exists_native_soft_parent_energy_bound};
the physical construction and minimum comparison are also exposed in
\leanid{TNLean.MPS.AreaLaw.ALVV.exists_soft_parent_structure}.

The norm statement concerns $R-e_{\mathrm s}I$, not $R$. A known rough
parent energy could instead be subtracted when constructing an operator;
the scalar offset is then bounded by the corresponding energy errors.
An unknown scalar normalization is not by itself an obstruction to a
normalized viable-space construction.

\section{The remaining energy-estimation argument}

At source line 1631, Figure \texttt{fig:algo1} uses the separate
untrimmed powered family image of the two child spaces as the energy
candidate, rather than the retained, trimmed output of the tree.
Lines 1793--1807 combine its viability for a generated soft target with
a norm estimate for the same generated operator. The estimate above is
on the original ground space. It must be transferred quantitatively to
that actual soft target before it can be applied to this candidate.
Neither that transfer, the recursive one-sided energy invariant, nor a
running-time bound is a conclusion of the present auxiliary theorem.

The false printed soft-energy and cluster estimates retain their
counterexamples and separate corrections; see
\href{https://sirui-lu.com/TNLean/paper-gaps/alvv16_soft_regional_residuals.pdf}{the regional-residual note}
and
\href{https://sirui-lu.com/TNLean/paper-gaps/alvv16_cluster_approximation.pdf}{the cluster-approximation note}.
No printed approximation estimate is assumed in the construction above.
\end{document}
